% The R chapter of the WFDB-x user's manual. It follows the sections of the Python chapter
% (python.tex), in the same order. Each \apientry of the section "The high-level API" (and
% of the section on the command-line tools) documents one function exported by r/NAMESPACE
% from r/R/highlevel.R, highlevel_annot.R, metadata.R, check.R and cli.R: its signature
% listing is the usage of the function, with the argument names and defaults of its
% definition in the code, and its params name every formal argument, in code order. The
% listings of style "example" were run with Rscript (the package installed from r/), in a
% copy of tests/data as the current directory, and print what their comments say.
\chapter{R}
\label{chap:r}
\chapterlanguage{r}

The R implementation of WFDB-x is the package \code{wfdbx}, in the directory \file{r} of
the repository. It is written in base R alone: it needs no other package, has no compiled
code, and never calls the C library. It reads and writes the same files as the original C
library, and its tools print the same output.

It offers WFDB-x in three ways:
\begin{itemize}
\item \textbf{The high-level API} (\cref{sec:r-api}): functions that read and write whole
  records, headers, annotation files and header metadata as R matrices, lists and data
  frames, in one call, with named arguments, and that stop with R errors on failure. They
  read signal files in large blocks with vectorized code, and are the fastest way to read
  a record. They are an addition of WFDB-x: the C package has no equivalent. They are the
  recommended interface for new code and analysis.
\item \textbf{The WFDB library functions} under their C names (\code{isigopen},
  \code{getvec}, \code{getann}, \code{putann}, \ldots), with the behavior of the C library
  (\cref{sec:r-library}). They are documented once, for all the languages, in
  \cref{chap:library}.
\item \textbf{The command-line tools} (\code{rdsamp}, \code{rdann}, \code{wrsamp},
  \code{xform} and 29 others), with the options, output and exit status of the C tools,
  run with \cmd{Rscript} through the driver script \file{wfdbx.R}, or from R with
  \code{wfdbx\_tool} (\cref{sec:r-tools}). They are documented in \cref{chap:tools}.
\end{itemize}
The high-level functions return exactly the samples and annotations that the library
functions (and the C library) give.

% ==========================================================================================
\section{Installation and requirements}
\label{sec:r-install}

\subsection{Requirements}

WFDB-x for R needs R 4.1 or later, and nothing else: the package uses base R only (its
one import, \code{tools}, is part of R). It has no compiled code and needs no C compiler,
no WFDB library and no libFLAC (the FLAC formats 508, 516 and 524 are read and written by
WFDB-x's own codec, in R). It runs on Linux and macOS. On Windows, the command-line tools
write their output as text, so NUL bytes in it are lost (\cref{sec:r-tools}).

Two programs are used if they are installed: \cmd{curl}, to download remote files (base
R's \code{url()} is used otherwise), and \cmd{stat}, to recognize an output file that is
also an input file under another path. The tests also use \cmd{python3}
(\cref{sec:r-testing}).

\subsection{Installing the package}

From the root of a copy of the repository, the package is in the directory \file{r}:

\begin{lstlisting}[style=shell]
R CMD INSTALL r                   # install the package wfdbx
\end{lstlisting}

or, from R, with the path of the directory \file{r}:

\begin{lstlisting}[language=R]
install.packages("r", repos = NULL, type = "source")
library(wfdbx)
\end{lstlisting}

The package is attached with \code{library(wfdbx)}. Its version is given by
\code{packageVersion("wfdbx")}, and by the driver script of the tools:

\begin{lstlisting}[style=shell]
Rscript -e 'cat(format(packageVersion("wfdbx")), "\n")'    # 1.2.1
Rscript r/inst/bin/wfdbx.R --version                        # WFDB-x 1.2.1
\end{lstlisting}

The help pages of the package (\code{?read\_record}, \code{?read\_metadata},
\code{?isigopen}, \ldots) document each function; \code{?wfdbx} gives an overview.

\subsection{The command-line tools}

The tools are run by the driver script \file{wfdbx.R} with \cmd{Rscript}. In a copy of the
repository, the script is \file{r/inst/bin/wfdbx.R}, and it runs the tools without the
package being installed (it then loads the sources of \file{r/R}); in an installed package,
its path is \code{system.file("bin", "wfdbx.R", package = "wfdbx")}
(\cref{sec:r-tools}).

\subsection{The sample record}
\label{sec:r-resources}

The repository includes the sample record \code{100s}, the sample record of the WFDB
Software Package: the first minute of record 100 of the MIT-BIH Arrhythmia Database, two
ECG signals (\code{MLII} and \code{V5}) sampled at 360\,Hz, with the reference beat
annotations (annotator \code{atr}). Its files (\file{100s.hea}, \file{100s.dat},
\file{100s.atr}) are in the directory \file{tests/data}, with the other test records of
WFDB-x; the R package does not install them. Work that uses the MIT-BIH Arrhythmia
Database cites it as described on PhysioNet.

% ==========================================================================================
\section{Quick start}
\label{sec:r-quickstart}

This example reads the header, ten seconds of the signals and the annotations of the
sample record \code{100s}, from the directory \file{tests/data} of a copy of the
repository (\code{input\_dir}; the path is that of the copy), and computes the mean heart
rate. It then writes the first ten seconds as a new record, \code{first10}, in the current
directory, with two metadata entries in its header.

\begin{lstlisting}[language=R,style=example]
library(wfdbx)

data <- "wfdb-x/tests/data"             # the directory of the sample record 100s

# The header only: no samples are read.
hdr <- read_header("100s", input_dir = data)
cat(hdr$nsig, hdr$fs, hdr$nsamples, "\n")           # 2 360 21600
print(sapply(hdr$siginfo, function(s) s$desc))     # "MLII" "V5"

# Seconds 10 to 20 of both signals, in physical units: one row per sample, one
# column per signal.
rec <- read_record("100s", sampfrom = "0:10", sampto = "0:20", physical = TRUE,
                   input_dir = data)
print(dim(rec$signals))                 # 3600 2
print(rec$signals[1, ])                 # -0.390 -0.275

# The beat annotations, and the mean heart rate.
ann <- read_annotations("100s", "atr", input_dir = data)
cat(nrow(ann), head(ann$symbol, 3), head(ann$sample, 3), "\n")   # 75 + N N 18 77 370
beats <- ann$sample[isqrs(ann$code) == 1]
rr <- diff(beats) / attr(ann, "fs")     # RR intervals, in seconds
cat(length(beats), "beats,", round(60 / mean(rr), 1), "beats per minute\n")   # 74 beats, 73.9 ...

# Write the first 10 seconds, in ADC units, as the record "first10" (first10.hea and
# first10.dat, in the current directory), with the signal specifications of 100s.
x <- read_record("100s", sampto = "0:10", input_dir = data)
write_record("first10", x$signals, fs = x$fs, fmt = 16, gains = x$gains,
             baselines = x$baselines,
             adczeros = sapply(x$siginfo, function(s) s$adczero),
             adcres = sapply(x$siginfo, function(s) s$adcres), units = "mV",
             names = x$names, metadata = list(Source = "MIT-BIH record 100"))

# Header metadata: "# Key: value" lines of the header.
set_metadata("first10", "Age", "69")
print(read_metadata("first10"))         # c(Source = "MIT-BIH record 100", Age = "69")
print(check_record("first10"))          # no problems found
\end{lstlisting}

The header \file{100s.hea} gives no units: the example writes \code{"mV"} explicitly,
since WFDB-x never assumes units, a gain, an ADC resolution or a sampling frequency that
a record does not give (\cref{chap:behavior}).

% ==========================================================================================
\section{The high-level API}
\label{sec:r-api}

The high-level API is made of the functions of this section and of the lists and data
frames that they return (\cref{sec:r-results}). They are defined in the files
\file{highlevel.R}, \file{highlevel\_annot.R}, \file{metadata.R} and \file{check.R} of
\file{r/R}, and documented by the help pages \code{?read\_record} (which also covers
\code{read\_header} and \code{write\_record}), \code{?read\_annotations} (also
\code{write\_annotations}), \code{?check\_record} and \code{?read\_metadata} (the five
metadata functions).

\subsection{Conventions}
\label{sec:r-conventions}

\paragraph{Record names and directories.} A record is named as in the other parts of
WFDB-x (\cref{chap:records,chap:files}), without the \code{.hea} extension (the readers
also accept a trailing \code{.hea}; \code{write\_record} does not). A record name is a
full path, absolute (\code{"/data/mitdb/100"}) or in the home directory
(\code{"\textasciitilde/Documents/100s"}), used as given; a relative path, such as \code{"100s"} or
\code{"data/100s"}, read from the input directory if one is set, otherwise from the
current directory, and written to the output directory if one is set, otherwise to the
current directory (missing directories are created); or a URL
(\code{"https://physionet.org/files/mitdb/1.0.0/100"}), read only if network access is
on, which cannot be written. A writer given a name with a path creates the files in that
directory, and the header names the record and its signal file without the path
(\code{rec}, \code{rec.dat}), so that the record opens under any path that leads to it;
in legacy mode, as in the original, the header gives the name with its path. The input
and output directories are the current directory unless they are set, with \code{wfdbx\_set\_input\_dir} and
\code{wfdbx\_set\_output\_dir} (\cref{sec:r-settings}) or with the environment variables
\code{WFDBX\_INPUT\_DIR} and \code{WFDBX\_OUTPUT\_DIR}. URLs are read only when network
access is on. The functions that read take the arguments \code{input\_dir} and
\code{network}, and those that write take \code{dir}: they apply to that one call, and the
previous settings are restored when it returns, also after an error. \code{NULL}, the
default of these arguments, keeps the current setting (for \code{dir}, \code{""} too).
\code{input\_dir} applies to relative names only, not to full paths or URLs. With
\code{dir}, the name must be a plain name (\code{"rec"}): the files are created in
\code{dir}, and the record opens from that directory under its plain name.
The signal files that a header names are read from the header's directory. A name with
another extension, such as \code{"rec.edf"}, is read as an EDF file.

\paragraph{Times.} The arguments \code{sampfrom} and \code{sampto} are sample numbers
(whole numbers) or WFDB time strings, as \code{strtim} reads them: \code{"0:10"} (10
seconds), \code{"1:02:03.5"}, \code{"s3600"} (sample 3600), \code{"e"} (the end of the
record), or a time of day in brackets, \code{"[14:00:00]"}. A time in seconds needs the
sampling frequency; when the record does not give it, such a time is an error.

\paragraph{Per-signal and per-annotation arguments.} The arguments that describe each
signal (\code{fmt}, \code{gains}, \code{units}, \ldots) or each annotation
(\code{subtype}, \code{chan}, \ldots) take one value for all, or a vector (or, for
\code{units} and \code{names}, a list) with exactly one value per signal or per
annotation; any other length is an error.

\paragraph{Numbers and text.} Numbers are doubles: samples, times and counts are
\code{double} vectors (exact for every 32-bit sample and for times up to $2^{53}$).
Signal numbers start at 0, as in WFDB. Text is handled as bytes: what the files hold is
returned as character strings of their bytes, never re-encoded, and character strings
given are written as their bytes.

\paragraph{Signal parameters.} No sampling frequency, gain, ADC resolution or units is
ever assumed. Where one is needed and the record does not give it, it must be given as an
argument; otherwise the function stops with an error. A record whose header gives no
sampling frequency has \code{fs = 0}.

\paragraph{Legacy mode.} With \code{legacy = TRUE}, a function behaves as in the original
WFDB package, with its assumed signal parameters and its search for files
(\cref{chap:behavior}); the library's legacy mode is on for the duration of the call
only. With \code{legacy = FALSE} (the default), the current setting of
\code{wfdbx\_set\_legacy} applies, which is off unless it was turned on.

\paragraph{State.} The functions that read and write records and annotation files close
anything opened with the WFDB library functions (they call \code{wfdbquit()}) before
they start, and close their own files before they return, whatever the outcome.
\code{read\_record} sets the getvec mode for its own reading only. \code{read\_record},
\code{read\_header}, \code{read\_annotations}, \code{write\_record},
\code{write\_annotations} and \code{check\_record} leave the getvec mode
(\code{getgvmode}) and the time resolution of output annotators (\code{getafreq}) as they
were, also after an error: opening an EDF record sets both, as in C (\code{WFDB\_HIGHRES}
is added to the getvec mode, and the time resolution becomes the record's sampling
frequency when none is set), and \code{wfdbquit()} resets neither; the \code{fs} of
\code{write\_annotations} applies to its file only. The metadata functions open no
signals and leave the library's state as it is.

\paragraph{The examples} of this section read the test records of WFDB-x (among them
\code{100s}) from the current directory, a copy of \file{tests/data}, and write records of
their own there. They were run with \cmd{Rscript}, and print what their comments say.

% ------------------------------------------------------------------------------------------
\subsection{Reading records}
\label{sec:r-reading}

\apientry{read\_record}{Read the signals of a record}
\begin{lstlisting}[language=R,style=signature]
read_record(record, channels = NULL, sampfrom = 0, sampto = NULL,
            physical = FALSE, highres = FALSE, pad = FALSE,
            input_dir = NULL, network = NULL,
            legacy = FALSE, checksum_warning = TRUE)
\end{lstlisting}
Reads the signals of a record, or part of it, into a numeric matrix with one row per
sample and one column per signal, together with the description of the record. The
samples are those that \code{getvec} returns (\cref{chap:library}), read in large blocks.
Invalid samples are \code{WFDB\_INVALID\_SAMPLE} ($-32768$) in ADC units, and \code{NA}
in physical units. When the read reaches the end of the record, the samples are compared
with the checksums of the header (\cref{sec:r-warnings}).
\begin{params}
\item[record] Record name (a character string), without the \code{.hea} extension (a
  trailing \code{".hea"} is also accepted): a full path, absolute
  (\code{"/data/mitdb/100"}) or in the home directory
  (\code{"\textasciitilde/Documents/100s"}), used as given; a relative path (\code{"100s"},
  \code{"data/100s"}), read from the input directory if one is set, otherwise from the
  current directory; or a URL (\code{"https://physionet.org/files/mitdb/1.0.0/100"}),
  read only if network access is on.
\item[channels] The signals to read, in the order in which they are returned: signal
  numbers (from 0) or signal descriptions (as in \code{names}, such as \code{"MLII"}).
  Default \code{NULL}: all the signals, in the order of the header.
\item[sampfrom] The first sample to read: a sample number (a whole number $\ge 0$) or a
  time string. Default \code{0}: the beginning of the record.
\item[sampto] The sample after the last one to read: a sample number (a whole number) or
  a time string. A value at or before \code{sampfrom} reads no samples. Default
  \code{NULL}: the end of the record.
\item[physical] \code{FALSE} (the default) to return ADC units; \code{TRUE} to return
  physical units, computed as $(\mathit{adc} - \mathit{baseline}) / \mathit{gain}$. Every
  signal read must then have a gain.
\item[highres] \code{FALSE} (the default) to read a multi-frequency record at its frame
  frequency, with the samples of the faster signals averaged to one per frame;
  \code{TRUE} to read every sample, at the highest sampling frequency of the record (the
  samples of the slower signals are repeated). \code{fs} of the result is then that
  frequency.
\item[pad] \code{TRUE} to replace each invalid sample by the previous valid sample of the
  same signal. Default \code{FALSE}.
\item[input\_dir] The input directory for this call only (optional): a local directory, a
  \code{\textasciitilde} path or a URL, under which relative record names are read. It does not
  apply to full paths or URLs. Default \code{NULL}: the current setting
  (\code{wfdbx\_set\_input\_dir}).
\item[network] \code{TRUE} to allow reading URLs for this call only, \code{FALSE} to forbid
  it. Default \code{NULL}: the current setting (\code{wfdbx\_set\_network}).
\item[legacy] \code{TRUE} to behave as the original WFDB package (legacy mode, for this
  call only). Default \code{FALSE}.
\item[checksum\_warning] \code{TRUE} (the default) to issue an R warning when the samples
  do not match the checksums of the header; \code{FALSE} for no warning. The samples are
  returned either way, and \code{checksum\_ok} of the result is \code{FALSE}.
\end{params}
\returns{A list (\cref{tab:r-record}).}
\raises{Stops with an error if the header or a signal file cannot be read (``can't read
record REC''), if the record has no signals, if \code{sampfrom} cannot be reached (``can't
seek to sample N of record REC''), or if a signal file is shorter than its header says
(``can't read signal data of record REC (getvec status -3)''); if \code{sampfrom} is
negative (``sampfrom must not be negative''), or \code{sampfrom} or \code{sampto} is not a
whole number (a fraction, \code{NA}, \code{NaN} or an infinite value: ``sampfrom must be
a whole number''); if a channel is not in the record (``record REC has no signal C''), if
a time string needs the sampling frequency and the header gives none (``sampling
frequency of record REC is not specified''), or if \code{physical = TRUE} and a signal
read has no gain (``signal K of record REC is uncalibrated''). Warns on a checksum error
(see \code{checksum\_warning}).}
\begin{lstlisting}[language=R,style=example]
library(wfdbx)

# Seconds 10 to 20 of lead MLII, in physical units.
rec <- read_record("100s", channels = "MLII", sampfrom = "0:10", sampto = "0:20",
                   physical = TRUE)
print(dim(rec$signals))                  # 3600 1
cat(rec$names, rec$fs, rec$sampfrom, "\n")   # MLII 360 3600

# The first 3 samples of both signals, in ADC units, V5 first.
rec <- read_record("100s", channels = c(1, 0), sampto = 3)
print(rec$signals)                       # three rows of 1011 995
print(rec$gains)                         # 200 200
print(rec$baselines)                     # 1024 1024
str(rec$units)                           # a list of two NULL: no units in the header

# The whole record: the checksums are verified.
rec <- read_record("100s")
cat(dim(rec$signals), rec$checksum_ok, "\n")   # 21600 2 TRUE
\end{lstlisting}
A multi-frequency record (here the test record \code{mfreq}, whose signals have 1, 2
and 4 samples per frame at a frame frequency of 100\,Hz) is read at its frame frequency,
or at its highest sampling frequency:
\begin{lstlisting}[language=R,style=example]
library(wfdbx)

rec <- read_record("mfreq")
cat(rec$fs, dim(rec$signals), "\n")      # 100 600 3
rec <- read_record("mfreq", highres = TRUE)
cat(rec$fs, dim(rec$signals), "\n")      # 400 2400 3
print(sapply(rec$siginfo, function(s) s$spf))   # 1 2 4
\end{lstlisting}

\apientry{read\_header}{Read the header of a record}
\begin{lstlisting}[language=R,style=signature]
read_header(record, input_dir = NULL, network = NULL,
            legacy = FALSE)
\end{lstlisting}
Reads the header of a record, without opening its signal files. The result describes the
record and each of its signals, and holds the info strings and the metadata entries of the
header. For a multi-segment record, the master header is read, and \code{nsamples} is the
length of the whole record. The metadata are read from the \file{.hea} file only; an EDF
file, which has no \file{.hea} file, and a header read from the standard input
(\code{"-"}) have none.
\begin{params}
\item[record] Record name (a character string), without the \code{.hea} extension (a
  trailing \code{".hea"} is also accepted): a full path, absolute
  (\code{"/data/mitdb/100"}) or in the home directory
  (\code{"\textasciitilde/Documents/100s"}), used as given; a relative path (\code{"100s"},
  \code{"data/100s"}), read from the input directory if one is set, otherwise from the
  current directory; or a URL (\code{"https://physionet.org/files/mitdb/1.0.0/100"}),
  read only if network access is on.
\item[input\_dir] The input directory for this call only (optional): a local directory, a
  \code{\textasciitilde} path or a URL, under which relative record names are read. It does not
  apply to full paths or URLs. Default \code{NULL}: the current setting
  (\code{wfdbx\_set\_input\_dir}).
\item[network] \code{TRUE} to allow reading URLs for this call only, \code{FALSE} to
  forbid it. Default \code{NULL}: the current setting.
\item[legacy] \code{TRUE} to read the header as the original WFDB package does. Default
  \code{FALSE}.
\end{params}
\returns{A list of class \code{"wfdbx\_header"} (\cref{tab:r-header}).}
\raises{Stops with an error if the header cannot be found or read (``can't read header
for record REC''). Warns, as \code{read\_metadata} does, for each key of which the header
has more than one entry.}
\begin{lstlisting}[language=R,style=example]
library(wfdbx)

hdr <- read_header("100s")
cat(hdr$name, hdr$fs, hdr$nsig, hdr$nsamples, "\n")   # 100s 360 2 21600
for (s in hdr$siginfo)
  cat(s$desc, s$fname, s$fmt, s$gain, s$adcres, is.null(s$units), "\n")
# MLII 100s.dat 212 200 11 TRUE
# V5 100s.dat 212 200 11 TRUE
print(hdr$comments[1])          # " 69 M 1085 1629 x1": the first info string
print(is.null(hdr$base_time))   # TRUE: the header gives no start time
print(hdr$metadata)
# c("Produced by xform from record 100, beginning at 0" = "0"): the last info string
# contains a colon, so it is an entry
\end{lstlisting}

\apientry{read\_annotations}{Read an annotation file}
\begin{lstlisting}[language=R,style=signature]
read_annotations(record, annotator, sampfrom = 0, sampto = NULL, fs = NULL,
                 input_dir = NULL, network = NULL,
                 legacy = FALSE)
\end{lstlisting}
Reads the annotations of an annotation file, in file order, into a data frame with one
row per annotation. The sampling frequency, needed to convert times in seconds, is that of
the record's header, or, for a record without a header or whose header gives none, the
argument \code{fs}.
\begin{params}
\item[record] Record name (a character string), without the \code{.hea} extension (a
  trailing \code{".hea"} is also accepted): a full path, absolute
  (\code{"/data/mitdb/100"}) or in the home directory
  (\code{"\textasciitilde/Documents/100s"}), used as given; a relative path (\code{"100s"},
  \code{"data/100s"}), read from the input directory if one is set, otherwise from the
  current directory; or a URL (\code{"https://physionet.org/files/mitdb/1.0.0/100"}),
  read only if network access is on. Its header, if there is one, gives
  the sampling frequency.
\item[annotator] Annotator name: the suffix of the annotation file. \code{"atr"} reads the
  file \file{RECORD.atr}.
\item[sampfrom] The first time to read: a sample number or a time string. Default
  \code{0}: the beginning of the file.
\item[sampto] The time at which to stop: annotations at or after it are not read. A
  sample number or a time string. Default \code{NULL}: the end of the file.
\item[fs] The sampling frequency, in Hz (a number $> 0$), to use when the record has no
  header or its header gives none; it is ignored otherwise. Default \code{NULL}: none.
\item[input\_dir] The input directory for this call only (optional): a local directory, a
  \code{\textasciitilde} path or a URL, under which relative record names are read. It does not
  apply to full paths or URLs. Default \code{NULL}: the current setting
  (\code{wfdbx\_set\_input\_dir}).
\item[network] \code{TRUE} to allow reading URLs for this call only, \code{FALSE} to
  forbid it. Default \code{NULL}: the current setting.
\item[legacy] \code{TRUE} to behave as the original WFDB package (which assumes 250\,Hz
  for a record without a frequency). Default \code{FALSE}.
\end{params}
\returns{A data frame (\cref{tab:r-annotations}), with the attributes \code{record},
\code{annotator} and \code{fs}. \code{attr(ann, "fs")} is the frequency in which the times
are counted: the one written in the annotation file, if there is one, or else the
sampling frequency; it is 0 if none of them is known.}
\raises{Stops with an error if the annotation file cannot be opened (``can't read
annotator ANN for record REC'') or read, or if it ends before its end-of-file marker
(``can't read annotation file ANN of record REC (getann status -3)''); if \code{fs} is
given and is not $> 0$ (``fs must be > 0''); or if a time string needs the sampling
frequency and none is known.}
\begin{lstlisting}[language=R,style=example]
library(wfdbx)

ann <- read_annotations("100s", "atr", sampfrom = "0:10", sampto = "0:30")
cat(nrow(ann), attr(ann, "fs"), "\n")    # 24 360
print(head(ann, 3))
#   sample symbol code subtype chan num  aux
# 1   3862      N    1       0    0   0 <NA>
# 2   4170      N    1       0    0   0 <NA>
# 3   4466      N    1       0    0   0 <NA>

# The whole file: the first annotation is a rhythm change, with an aux string.
ann <- read_annotations("100s", "atr")
cat(ann$symbol[1], ann$code[1] == RHYTHM, ann$aux[1], "\n")   # + TRUE (N
print(sort(unique(ann$symbol)))          # "+" "A" "N"
\end{lstlisting}
For a record without a header, \code{fs} gives the frequency in which times in seconds
are converted (the library reports on the standard error that it cannot open the
header):
\begin{lstlisting}[language=R,style=example]
library(wfdbx)

write_annotations("beats", "qrs", c(100, 460, 820), symbol = "N")   # no header
ann <- read_annotations("beats", "qrs", sampto = "0:02", fs = 360)
cat(ann$sample, attr(ann, "fs"), "\n")   # 100 460 360
\end{lstlisting}

% ------------------------------------------------------------------------------------------
\subsection{Writing records}
\label{sec:r-writing}

\apientry{write\_record}{Write a header and a signal file}
\begin{lstlisting}[language=R,style=signature]
write_record(record, signals, fs, fmt = 16, gains, baselines = NULL,
             adczeros = 0, adcres = NULL, units = NULL, names = NULL,
             base_time = NULL, comments = NULL, physical = FALSE,
             samps_per_frame = NULL, dir = NULL,
             legacy = FALSE, metadata = NULL)
\end{lstlisting}
Writes signals as a new record: the header \file{RECORD.hea} and one signal file,
\file{RECORD.dat}, which holds all the signals (so they share one format). Existing files
of these names are overwritten. The sampling frequency, the gain and ADC resolution of
each signal, and the units of each calibrated signal must be given: there are no
defaults. The per-signal arguments (\code{fmt}, \code{gains}, \code{baselines},
\code{adczeros}, \code{adcres}, \code{units}, \code{names}) take one value for all the
signals, or a vector (or list) with one value per signal. All the arguments, the samples
included, are checked before any file is created; whatever happens, the output signals
are closed (\code{wfdbquit}) before the function returns.
\begin{params}
\item[record] Record name (a non-empty character string), without the \code{.hea}
  extension: a plain name (\code{"rec"}) or a relative path (\code{"data/rec"}), created
  in the output directory if one is set, otherwise in the current directory, missing
  directories included; or a full path, absolute or in the home directory
  (\code{"\textasciitilde/Documents/rec"}), used as given. With a path, the files are created in that
  directory, and the header names the record and its signal file without the path
  (\code{rec}, \code{rec.dat}), so that the record opens under any path that leads to it
  (in legacy mode, as in the original, the header gives the name with its path). A URL
  cannot be written.
\item[signals] The samples: a numeric matrix with one row per sample and one column per
  signal; a vector is one signal. In ADC units (the default), the values are integers
  within the range of a 32-bit integer, and \code{NA} is written as an invalid sample;
  with \code{physical = TRUE}, they are in physical units.
\item[fs] The sampling frequency, in Hz: a finite number $> 0$. Required.
\item[fmt] The storage format, such as \code{16}, \code{212}, \code{24} or \code{516}
  (FLAC); one value, or one per signal, all equal (\cref{chap:formats}). Default
  \code{16}.
\item[gains] The gain of each signal, in ADC units per physical unit: a finite number;
  \code{0} marks an uncalibrated signal, which is allowed only for samples given in ADC
  units. Negative gains are allowed. Required (no default).
\item[baselines] The ADC value that corresponds to 0 physical units. Default
  \code{NULL}: the ADC zero of the signal.
\item[adczeros] The ADC value at the middle of the ADC range. Default \code{0}.
\item[adcres] The ADC resolution, in bits (an integer $\ge 1$). Required (default
  \code{NULL}: not given).
\item[units] The physical units, such as \code{"mV"}, as a character vector or a list.
  Required for each signal whose gain is not 0. Default \code{NULL}: not given.
\item[names] The signal descriptions, such as \code{"MLII"}, as a character vector or a
  list. Default \code{NULL}: no descriptions.
\item[base\_time] The start time, with an optional date: \code{"hh:mm:ss"} or
  \code{"hh:mm:ss dd/mm/yyyy"}, with optional decimals of the seconds. Default
  \code{NULL}: no start time.
\item[comments] A character vector of info strings, written to the header as comment
  lines (\code{"\#"} followed by the string). Default \code{NULL}: none.
\item[physical] \code{TRUE} if \code{signals} are in physical units. They are converted
  to ADC units as $\mathrm{round}(x \cdot \mathit{gain}) + \mathit{baseline}$ (halves
  rounded away from zero); \code{NA} becomes an invalid sample, and infinite values and
  values beyond the range of 32-bit ADC units are errors. Default \code{FALSE}: ADC units.
\item[samps\_per\_frame] For a multi-frequency record, the number of samples per frame of
  each signal (integers $\ge 1$). \code{signals} then has one row per frame and
  \code{sum(samps\_per\_frame)} columns: the samples of signal 0, then those of signal 1,
  and so on. \code{fs} is then the frame frequency. Default \code{NULL}: one sample per
  frame for every signal.
\item[dir] The directory (which must exist) in which to create the files, for this call
  only; \code{record} must then be a plain name, without a directory, and the record
  opens from that directory under its plain name. Default \code{NULL} (or \code{""}): the
  output directory.
\item[legacy] \code{TRUE} to behave as the original WFDB package: the gain defaults to
  200, and the ADC resolution to one based on the format. Default \code{FALSE}.
\item[metadata] Metadata entries (\cref{chap:metadata}) to write to the header after the
  comments, as \code{\# Key: value} (through \code{putinfo}): a named list or a named
  character vector, such as \code{list(Age = "63")} or the result of
  \code{read\_metadata}. Keys and values follow the rules of \code{set\_metadata}; a
  \code{NULL} or \code{NA} value is an error, no two keys may differ only in letter case,
  and none may be the key of an entry in \code{comments}. Default \code{NULL}: none.
\end{params}
\returns{\code{NULL}, invisibly.}
\raises{Stops with an error if \code{record} is empty (``record name must not be
empty''); if \code{signals} has no signals (``at least one signal is required''); if
\code{fs}, a gain, an ADC resolution or the units of a calibrated signal is missing or
invalid (``fs must be finite and > 0'', ``gains must be specified'', ``gains must be
finite'', ``ADC resolution(s) must be specified'', ``units must be specified''); if
\code{base\_time} cannot be parsed (``invalid base\_time 'TEXT'\,''); if a vector of
per-signal values does not have one value per signal (``NAME must have one value per
signal''); if a signal in physical units has a gain of 0 (``signal K is uncalibrated'');
if a sample is not valid (``samples in ADC units must be integers'', ``samples in ADC
units must be within the range of a 32-bit integer''); if the signals do not all have the
same format; if a \code{samps\_per\_frame} value is not an integer $\ge 1$ (``samps\_per\_frame
must hold integers >= 1''), or the number of columns of \code{signals} is not their sum
(``signals has N columns, but samps\_per\_frame adds up to M''); if \code{metadata} is not
valid (an error of class \code{wfdbx\_argument\_error}, \cref{sec:r-warnings}); if
\code{record} contains a directory and \code{dir} is given, or \code{dir} does not exist
(``directory DIR does not exist''); and if the signal file or the header cannot be created,
or cannot be written completely, for example on a full disk (``write error on record
REC'').}
\begin{lstlisting}[language=R,style=example]
library(wfdbx)

# Two signals of random numbers, in millivolts, at 360 Hz: the record "rand"
# (rand.hea and rand.dat).
set.seed(1)
x <- matrix(rnorm(7200), ncol = 2)
write_record("rand", x, fs = 360, fmt = 16, physical = TRUE, gains = 1000, adcres = 16,
             units = "mV", names = c("A", "B"), base_time = "14:00:03 21/06/2024",
             comments = "random numbers", metadata = list(Generator = "rnorm"))
rec <- read_record("rand", physical = TRUE)
cat(dim(rec$signals), rec$base_time, "\n")      # 3600 2 [14:00:03 21/06/2024]
print(max(abs(rec$signals - x)) <= 0.0005)      # TRUE: rounded to 1/1000 mV
print(read_header("rand")$comments)             # "random numbers" " Generator: rnorm"

# An uncalibrated signal (gain 0, no units), in ADC units, in the directory "out".
dir.create("out", showWarnings = FALSE)
write_record("counts", 0:999, fs = 100, gains = 0, adcres = 12, dir = "out")
print(read_record("out/counts")$signals[1:3, 1])   # 0 1 2
\end{lstlisting}
A multi-frequency record: signal \code{fast} has 2 samples per frame, signal
\code{slow} one; each row of \code{signals} is one frame, with the two samples of
\code{fast} and then the sample of \code{slow}:
\begin{lstlisting}[language=R,style=example]
library(wfdbx)

frames <- 100                                   # 1 second at 100 frames per second
fast <- matrix(0:(2 * frames - 1), ncol = 2, byrow = TRUE)   # 200 samples per second
slow <- matrix(500, frames, 1)                  # 100 samples per second
write_record("mf", cbind(fast, slow), fs = 100, gains = 200, adcres = 12, units = "mV",
             names = c("fast", "slow"), samps_per_frame = c(2, 1))
rec <- read_record("mf", highres = TRUE)
cat(rec$fs, "\n")                               # 200
print(rec$signals[1:3, ])                       # rows 0 500, 1 500, 2 500
\end{lstlisting}

\apientry{write\_annotations}{Write an annotation file}
\begin{lstlisting}[language=R,style=signature]
write_annotations(record, annotator, sample, symbol = NULL, code = NULL,
                  subtype = 0, chan = 0, num = 0, aux = NULL, fs = NULL,
                  dir = NULL, legacy = FALSE, sort = NULL)
\end{lstlisting}
Writes an annotation file, \file{RECORD.ANNOTATOR}, in the MIT format; an existing file of
that name is overwritten. The record need not have a header. The annotation types are
given as \code{symbol} or as \code{code} (\cref{chap:formats}). Each per-annotation
argument is one value for all the annotations, or a vector with exactly one value per
annotation. If the annotations are not in time order, the file is sorted when it is
written, unless \code{sort = FALSE}. All the arguments are checked before the file is
created; whatever happens, the annotator is closed (\code{wfdbquit}) before the function
returns.
\begin{params}
\item[record] Record name (a character string), without the \code{.hea} extension (a
  trailing \code{".hea"} is also accepted): a plain name or a relative path
  (\code{"data/100s"}) is created in the output directory if one is set, otherwise in the
  current directory; a full path, absolute or in the home directory
  (\code{"\textasciitilde/Documents/100s"}), is used as given. The file is \file{RECORD.ANNOTATOR} at
  that place (\file{data/100s.qrs}). A URL cannot be written.
\item[annotator] Annotator name: the suffix of the file, made of letters, digits,
  \code{\_}, \code{-} and \code{\textasciitilde}.
\item[sample] The annotation times, in samples: integers within the range of a 64-bit
  integer, one per annotation. It may also be a data frame returned by
  \code{read\_annotations}, whose columns then give the other fields (\code{code},
  \code{subtype}, \code{chan}, \code{num} and \code{aux}) unless they are given.
\item[symbol] The mnemonics of the annotation types, such as \code{"N"} or \code{"V"}:
  one per annotation, or one string for all. An unknown mnemonic becomes code 0
  (\code{NOTQRS}). Default \code{NULL}: the types are given by \code{code}.
\item[code] The annotation codes (\code{NORMAL}, \code{PVC}, \ldots), instead of
  \code{symbol}: one per annotation, or one value for all. If both are given (as
  \code{read\_annotations} returns them), each symbol must be the mnemonic of its code (an
  empty symbol agrees with any code). Default \code{NULL}.
\item[subtype] The subtype of each annotation, stored as a signed byte (other values wrap
  around to $-128 \ldots 127$, as in C). Default \code{0}.
\item[chan] The channel number of each annotation, stored as an unsigned byte (other
  values wrap around to $0 \ldots 255$). Default \code{0}.
\item[num] The annotator number of each annotation, stored as a signed byte. Default
  \code{0}.
\item[aux] The auxiliary text of each annotation, with \code{NA} where there is none: a
  character vector, or one text for all. Texts longer than 255 bytes are cut. Default
  \code{NULL}: none.
\item[fs] If given (a finite number $> 0$), the frequency in which the times are
  counted, written to the file. Default \code{NULL}: none is written.
\item[dir] The directory (which must exist) in which to create the file, for this call
  only; \code{record} must then be a plain name, without a directory: the file is
  \file{DIR/RECORD.ANNOTATOR}. Default \code{NULL} (or \code{""}): the output directory.
\item[legacy] \code{TRUE} to behave as the original WFDB package (an out-of-order file is
  then sorted by running the \code{sortann} tool). Default \code{FALSE}.
\item[sort] \code{TRUE} to always sort the file, \code{FALSE} to keep the given order
  (without the library's advice to run \code{sortann}), for this call only. Default
  \code{NULL}: the environment variable \code{WFDBANNSORT} decides (sorting is on unless
  it is \code{0}).
\end{params}
\returns{The status of the last \code{putann} (0), invisibly.}
\raises{Stops with an error if neither \code{symbol} nor \code{code} is given (``symbol
or code must be given''); if both are given and disagree (``symbols and codes disagree at
annotation K'', K counting from 0); if a vector of per-annotation values does not have one
value per annotation (``NAME must have one value per annotation''); if a time is not an
integer within the range of a 64-bit integer; if a code, subtype, channel or annotator
number is \code{NA}, \code{NaN} or infinite (``NAME must not be NA, NaN or infinite''); if
\code{fs} is not finite and $> 0$; if \code{sort} is not \code{NULL}, \code{TRUE} or
\code{FALSE}; if \code{record} contains a directory and \code{dir} is given, or \code{dir}
does not exist; and if the annotation file cannot be created (``can't create annotator ANN
for record REC'') or cannot be written completely (``write error on annotator ANN for
record REC'').}
\begin{lstlisting}[language=R,style=example]
library(wfdbx)

# Three beats and a rhythm change, by symbol, with an aux string for the rhythm.
write_annotations("100s", "test", c(18, 77, 370, 662),
                  symbol = c("+", "N", "V", "N"), aux = c("(N", NA, NA, NA))
ann <- read_annotations("100s", "test")
print(ann$symbol)                        # "+" "N" "V" "N"
print(ann$aux)                           # "(N" NA NA NA

# A copy of the reference annotations: the data frame gives all the fields.
a <- read_annotations("100s", "atr")
write_annotations("100s", "copy", a)
b <- read_annotations("100s", "copy")
print(identical(b$sample, a$sample) && identical(b$symbol, a$symbol) &&
      identical(b$aux, a$aux))           # TRUE

# Annotations out of time order are sorted when the file is written (the library
# reports it on the standard error).
write_annotations("100s", "unsorted", c(500, 100), code = NORMAL)
print(read_annotations("100s", "unsorted")$sample)   # 100 500
\end{lstlisting}

% ------------------------------------------------------------------------------------------
\subsection{Checking records}
\label{sec:r-checking}

\apientry{check\_record}{Check a record and its annotation files}
\begin{lstlisting}[language=R,style=signature]
check_record(record, annotators = character(0), input_dir = NULL,
             network = NULL)
\end{lstlisting}
Checks the header of a record, its signal specifications, every sample of its signal
files, each segment of a multi-segment record, and the annotation files named, and
reports each problem found. A problem in the record never stops with an error: it is
reported, with a severity (\code{"error"}: the record is damaged or does not follow the
WFDB format, and programs will read wrong data or fail; \code{"warning"}: the record can
be read, but something is missing or unusual, and may give wrong results; \code{"info"}:
a fact worth knowing, such as invalid samples). The problems, their codes and their
messages are the same in the seven languages (\cref{tab:python-issue-codes}). The checks
run in normal mode, even if legacy mode is on, and the library prints no messages during
the call; both settings are restored afterwards. They are similar to those of Benjamin
Moody's \code{wfdbcheck}.
\begin{params}
\item[record] Record name (a character string), without the \code{.hea} extension (a
  trailing \code{".hea"} is also accepted): a full path, absolute
  (\code{"/data/mitdb/100"}) or in the home directory
  (\code{"\textasciitilde/Documents/100s"}), used as given; a relative path (\code{"100s"},
  \code{"data/100s"}), read from the input directory if one is set, otherwise from the
  current directory; or a URL (\code{"https://physionet.org/files/mitdb/1.0.0/100"}),
  read only if network access is on.
\item[annotators] A character vector of the annotator names whose files are checked, such
  as \code{"atr"}. Default \code{character(0)}: none.
\item[input\_dir] The input directory for this call only (optional): a local directory, a
  \code{\textasciitilde} path or a URL, under which relative record names are read. It does not
  apply to full paths or URLs. Default \code{NULL}: the current setting
  (\code{wfdbx\_set\_input\_dir}).
\item[network] \code{TRUE} to allow reading URLs for this call only, \code{FALSE} to
  forbid it. Default \code{NULL}: the current setting.
\end{params}
\returns{A list of class \code{"wfdbx\_check\_report"} (\cref{tab:r-checkreport}):
\code{report\$ok} is \code{TRUE} if no problem is an error, \code{report\$issues} is a data
frame with one row per problem, \code{format(report)} gives the text form, one line per
problem, and \code{print(report)} prints it (\code{no problems found} if there is none).}
\raises{Stops with an error if \code{record} is not a single string (``record must be a
string'') or is empty (``record must not be empty''), or if \code{annotators} is not a
character vector without \code{NA} (``annotator names must be strings''). A record that
cannot be read is reported (code \code{header\_unreadable}), not raised.}
\begin{lstlisting}[language=R,style=example]
library(wfdbx)

report <- check_record("100s", "atr")
print(report$ok)          # TRUE: no errors
report
# warning: signal 0: the units are not specified
# warning: signal 1: the units are not specified

# A damaged record: one sample of a 12-bit signal is overwritten with 2048.
write_record("damaged", 0:99, fs = 100, gains = 100, adcres = 12, units = "mV",
             names = "ramp")
con <- file("damaged.dat", "r+b")
invisible(seek(con, 20, rw = "write"))  # sample 10 (2 bytes per sample in format 16)
writeBin(as.raw(c(0x00, 0x08)), con)    # 0x0800 = 2048, little-endian
close(con)
report <- check_record("damaged")
print(report$ok)          # FALSE
print(report$issues[, c("severity", "code", "where")])
#   severity         code    where
# 1    error sample_range signal 0
# 2    error     checksum signal 0
print(format(report)[1])
# "error: signal 0: 1 sample is outside the ADC range (-2048 to 2047): sample 10 (value 2048)"
\end{lstlisting}

% ------------------------------------------------------------------------------------------
\subsection{Header metadata}
\label{sec:r-metadata}

These functions read and change the metadata of a header: its info strings of the form
\code{\# Key: value}, such as \code{\# Age: 63} (\cref{chap:metadata}). In short: keys are
compared without regard to letter case (A-Z and a-z only); a key must not be empty and
must not contain \code{:}, \code{\#}, a line break or a NUL byte, and a value must not
contain a line break or a NUL byte; spaces at either end of keys and values are removed.
Keys and values are handled as bytes: a character string is taken as its bytes, never
re-encoded, and a raw vector is also accepted, as given. A change rewrites the header in
one step (a temporary file in the same directory, renamed over the header, keeping its
permission bits), keeping every other byte of it; all the arguments are checked first, so
an error leaves the header unchanged, and nothing is written if nothing changes (unless
\code{dir} is given). Only the \file{.hea} file is read (not a \file{.info} file); for a
multi-segment record, it is the master header, unless a segment is named. The output
directory is not used: a change is written to the header that was read, or to \code{dir}.

The entries are a named character vector: the names are the keys and the elements the
values, in file order, with the spelling of the file, such as \code{c(Age = "63", Sex =
"F")}. Since reading returns each key once, \code{md[["Age"]]} gives a value (with the
file's spelling of the key; \code{get\_metadata} ignores letter case). The same form, or
a named list, is what \code{update\_metadata} and the argument \code{metadata} of
\code{write\_record} take, so the entries of one header can be written to another.

Metadata often holds information about patients: WFDB-x stores it as plain text, exactly
as given, and does not de-identify it.

\begin{phibox}
A header can hold identifying data about a patient (a name, a date of birth, a medical
record number). The metadata functions read and write such data as given. Records that
are shared must be de-identified first.
\end{phibox}

\apientry{read\_metadata}{Read the metadata of a record}
\begin{lstlisting}[language=R,style=signature]
read_metadata(record, input_dir = NULL, network = NULL)
\end{lstlisting}
Reads the metadata entries of the header of a record, in file order. An info string is an
entry if it contains a colon: the key is the text before the first colon, the value the
rest of the line, both without spaces at either end. Info strings without a colon are free
comments, which are not returned.
\begin{params}
\item[record] Record name (a character string), without the \code{.hea} extension (a
  trailing \code{".hea"} is also accepted): a full path, absolute
  (\code{"/data/mitdb/100"}) or in the home directory
  (\code{"\textasciitilde/Documents/100s"}), used as given; a relative path (\code{"100s"},
  \code{"data/100s"}), read from the input directory if one is set, otherwise from the
  current directory; or a URL (\code{"https://physionet.org/files/mitdb/1.0.0/100"}),
  read only if network access is on.
\item[input\_dir] The input directory for this call only (optional): a local directory, a
  \code{\textasciitilde} path or a URL, under which relative record names are read. It does not
  apply to full paths or URLs. Default \code{NULL}: the current setting
  (\code{wfdbx\_set\_input\_dir}).
\item[network] \code{TRUE} to allow reading URLs for this call only, \code{FALSE} to
  forbid it. Default \code{NULL}: the current setting.
\end{params}
\returns{A named character vector of the entries, in file order (an empty named character
vector if there are none). A key that appears more than once (ignoring letter case) is
returned once, with its first value, and an R warning is issued (\cref{sec:r-warnings}).}
\raises{An error of class \code{wfdbx\_file\_error} if the header cannot be found or read
(``can't read header for record REC''); of class \code{wfdbx\_argument\_error} if
\code{record} is empty or \code{"-"} (the standard input).}
\begin{lstlisting}[language=R,style=example]
library(wfdbx)

write_record("pat", rep(0, 360), fs = 360, gains = 200, adcres = 12, units = "mV",
             names = "ECG", comments = "recorded at rest",
             metadata = list(Age = "63", Sex = "F"))
print(read_metadata("pat"))             # c(Age = "63", Sex = "F")
print(read_header("pat")$comments)
# "recorded at rest" " Age: 63" " Sex: F": the free comment is not an entry
\end{lstlisting}

\apientry{get\_metadata}{Return the value of a metadata entry}
\begin{lstlisting}[language=R,style=signature]
get_metadata(record, key, input_dir = NULL, network = NULL)
\end{lstlisting}
Returns the value of one metadata entry of the header of a record.
\begin{params}
\item[record] Record name (as for \code{read\_metadata}).
\item[key] The key, such as \code{"Age"}: a character string or a raw vector. Letter case
  is ignored, and spaces at either end are removed.
\item[input\_dir] The input directory for this call only. Default \code{NULL}: the
  current setting.
\item[network] \code{TRUE} to allow reading URLs for this call only, \code{FALSE} to
  forbid it. Default \code{NULL}: the current setting.
\end{params}
\returns{The value, a character string (\code{""} for an entry with an empty value), or
\code{NULL} if the header has no entry for \code{key}. If it has more than one, the first
is returned, and an R warning is issued.}
\raises{An error of class \code{wfdbx\_file\_error} if the header cannot be found or
read; of class \code{wfdbx\_argument\_error} if \code{record} is empty or \code{"-"}, or
if \code{key} is not a valid key or not a string.}
\begin{lstlisting}[language=R,style=example]
library(wfdbx)

write_record("pat", rep(0, 360), fs = 360, gains = 200, adcres = 12, units = "mV",
             metadata = c(Age = "63", Sex = "F"))
print(get_metadata("pat", "age"))       # "63"
print(get_metadata("pat", "Height"))    # NULL
\end{lstlisting}

\apientry{set\_metadata}{Set the value of a metadata entry}
\begin{lstlisting}[language=R,style=signature]
set_metadata(record, key, value, input_dir = NULL, network = NULL, dir = NULL)
\end{lstlisting}
Sets the value of a metadata entry in the header of a record. If the header has an entry
for \code{key} (ignoring letter case), its line is replaced by \code{\# Key: value}, at
the same position, with the spelling of \code{key} given here; otherwise the entry is
added after the last info string of the header (or, if it has none, after its last line
that is not blank). The entry is written as \code{\# Key: value}, or \code{\# Key:} if
the value is empty. Every other line is kept, byte for byte, including its line ending.
\begin{params}
\item[record] Record name (as for \code{read\_metadata}). The header is changed where it
  was read, unless \code{dir} is given; a header read from a URL can only be written
  elsewhere, with \code{dir}.
\item[key] The key, such as \code{"Age"} or \code{"Date of birth"}: a character string
  or a raw vector. Spaces at either end are removed. It must not be empty, and must not
  contain \code{:}, \code{\#}, a line break or a NUL byte.
\item[value] The value: a character string, which may be a sentence, or a raw vector.
  Spaces at either end are removed. It must not contain a line break or a NUL byte; it
  may be empty.
\item[input\_dir] The input directory in which the header is read, for this call only.
  Default \code{NULL}: the current setting.
\item[network] \code{TRUE} to allow reading the header from a URL, for this call only,
  \code{FALSE} to forbid it. Default \code{NULL}: the current setting.
\item[dir] A directory, which must exist, in which to write the changed header, under the
  same file name, instead of changing the header itself; the signal files are not copied.
  The header is then written even if nothing changes. Default \code{NULL} (or \code{""}):
  the header itself is changed.
\end{params}
\returns{\code{NULL}, invisibly.}
\raises{An error of class \code{wfdbx\_file\_error} if the header cannot be read, if the
changed header cannot be written (``can't write header for record REC''), or if
\code{dir} does not exist (``directory DIR does not exist''); of class
\code{wfdbx\_argument\_error} if \code{key} or \code{value} is not valid or not a string,
if the header has more than one entry for \code{key} (``header of record REC has more than
one entry for key KEY (lines L1, L2); remove them first''), if it has no record line, if
it was read from a URL and \code{dir} is not given, or if \code{record} is empty or
\code{"-"}.}
\begin{lstlisting}[language=R,style=example]
library(wfdbx)

write_record("pat", rep(0, 360), fs = 360, gains = 200, adcres = 12, units = "mV",
             metadata = c(Age = "63", Sex = "F"))
set_metadata("pat", "AGE", "64")                                # replaces "# Age: 63"
set_metadata("pat", "Diagnosis", "sinus rhythm, no ectopy")     # a new entry
print(read_metadata("pat"))
# c(AGE = "64", Sex = "F", Diagnosis = "sinus rhythm, no ectopy")

# Write the changed header to another directory; pat.hea is left unchanged.
dir.create("changed", showWarnings = FALSE)
set_metadata("pat", "Sex", "M", dir = "changed")
cat(get_metadata("pat", "Sex"), get_metadata("changed/pat", "Sex"), "\n")   # F M
\end{lstlisting}

\apientry{remove\_metadata}{Remove a metadata entry}
\begin{lstlisting}[language=R,style=signature]
remove_metadata(record, key, input_dir = NULL, network = NULL, dir = NULL)
\end{lstlisting}
Removes every entry for \code{key} (ignoring letter case) from the header of a record, so
it also repairs a header that has more than one. Other lines are kept, byte for byte.
Nothing is written if the header has no entry for \code{key}, unless \code{dir} is given.
\begin{params}
\item[record] Record name (as for \code{set\_metadata}).
\item[key] The key (as for \code{set\_metadata}).
\item[input\_dir] The input directory in which the header is read, for this call only.
  Default \code{NULL}: the current setting.
\item[network] \code{TRUE} to allow reading the header from a URL, for this call only,
  \code{FALSE} to forbid it. Default \code{NULL}: the current setting.
\item[dir] A directory, which must exist, in which to write the changed header instead of
  changing the header itself. Default \code{NULL} (or \code{""}): the header itself is
  changed.
\end{params}
\returns{The number of lines removed (a number; 0 if the header has no entry for
\code{key}).}
\raises{As \code{set\_metadata}, except that more than one entry for \code{key} is not an
error.}
\begin{lstlisting}[language=R,style=example]
library(wfdbx)

write_record("pat", rep(0, 360), fs = 360, gains = 200, adcres = 12, units = "mV",
             metadata = c(Age = "63", Sex = "F"))
print(remove_metadata("pat", "sex"))    # 1
print(remove_metadata("pat", "Sex"))    # 0: there is no such entry now
print(read_metadata("pat"))             # c(Age = "63")
\end{lstlisting}

\apientry{update\_metadata}{Make several changes to the metadata}
\begin{lstlisting}[language=R,style=signature]
update_metadata(record, changes, input_dir = NULL, network = NULL, dir = NULL)
\end{lstlisting}
Makes several changes to the metadata of a record, in one write of the header. Each
change is applied in order, as \code{set\_metadata} does, or, for a value of \code{NULL}
or \code{NA}, as \code{remove\_metadata} does. All the changes are checked first: if one
is not valid, none is made.
\begin{params}
\item[record] Record name (as for \code{set\_metadata}).
\item[changes] A named list or a named character vector: the names are the keys, and each
  element a value (a character string or a raw vector), or \code{NULL} or \code{NA} to
  remove the entry, as in \code{list(Age = "64", Dx = NULL)} or \code{c(Age = "64", Dx =
  NA)}. No two keys may differ only in letter case. An empty list or vector makes no
  change.
\item[input\_dir] The input directory in which the header is read, for this call only.
  Default \code{NULL}: the current setting.
\item[network] \code{TRUE} to allow reading the header from a URL, for this call only,
  \code{FALSE} to forbid it. Default \code{NULL}: the current setting.
\item[dir] A directory, which must exist, in which to write the changed header instead of
  changing the header itself. Default \code{NULL} (or \code{""}): the header itself is
  changed.
\end{params}
\returns{\code{NULL}, invisibly.}
\raises{As \code{set\_metadata}; also an error of class \code{wfdbx\_argument\_error} if
two keys of \code{changes} differ only in letter case, or if \code{changes} is not a
named list or a named character vector (\code{NULL} included).}
\begin{lstlisting}[language=R,style=example]
library(wfdbx)

write_record("pat", rep(0, 360), fs = 360, gains = 200, adcres = 12, units = "mV",
             metadata = c(Age = "63", Sex = "F"))
update_metadata("pat", list(Age = "64", Sex = NULL, Dx = "164889003"))
print(read_metadata("pat"))             # c(Age = "64", Dx = "164889003")
tryCatch(update_metadata("pat", c(Dx = "426783006", dx = NA)),
         wfdbx_argument_error = function(e) print(conditionMessage(e)))
# "metadata key 'dx' is given more than once (keys are compared without regard to letter case)"
print(get_metadata("pat", "Dx"))        # "164889003": nothing was changed
\end{lstlisting}

% ------------------------------------------------------------------------------------------
\subsection{The results}
\label{sec:r-results}

The readers return plain R objects: \code{read\_record} a list
(\cref{tab:r-record}), \code{read\_header} a list of class \code{"wfdbx\_header"}
(\cref{tab:r-header}), \code{read\_annotations} a data frame (\cref{tab:r-annotations}),
\code{read\_metadata} a named character vector (\cref{sec:r-metadata}), and
\code{check\_record} a list of class \code{"wfdbx\_check\_report"}
(\cref{tab:r-checkreport}). The signals are described by \code{siginfo} lists
(\cref{tab:r-siginfo}), which are those of the WFDB library functions
(\cref{sec:r-library}). The elements of these lists are listed in the order in which they
appear.

\begin{xltabular}{\linewidth}{@{}lL@{}}
\caption{The elements of the list returned by \code{read\_record}}
\label{tab:r-record} \\
\toprule
\tabhead Element & \tabhead Contents \\
\midrule
\endfirsthead
\toprule
\tabhead Element & \tabhead Contents \\
\midrule
\endhead
\bottomrule
\endfoot
\code{name} & the record name, as given to \code{read\_record} \\
\code{fs} & the sampling frequency in Hz, or 0 if the header gives none; with
  \code{highres = TRUE}, the highest sampling frequency of a multi-frequency record \\
\code{signals} & a numeric matrix with one row per sample and one column per signal read:
  ADC units, or physical units with \code{physical = TRUE}; invalid samples are
  \code{WFDB\_INVALID\_SAMPLE}, or \code{NA} in physical units \\
\code{siginfo} & a list of \code{siginfo} lists, one per signal read, in the order of the
  columns; \code{desc} has no white space at either end, and \code{adcres} is 0 if the
  header does not give it \\
\code{names} & the description of each signal (a character vector), such as
  \code{"MLII"}, without the white space that a header may put at either end of it \\
\code{units} & the physical units of each signal, as a list, with \code{NULL} where the
  header gives none \\
\code{gains} & the gain of each signal, in ADC units per physical unit; 0 means that the
  signal is uncalibrated \\
\code{baselines} & the ADC value of each signal that corresponds to 0 physical units \\
\code{base\_time} & the start time and date as text, to the second, such as
  \code{"[14:00:03 21/06/2024]"} (\code{"[14:00:03]"} without a date), or \code{NULL} if
  the header gives none \\
\code{base\_time\_ms} & the start time in milliseconds after midnight, such as
  \code{50403500} for 14:00:03.500, or \code{NULL} if \code{base\_time} is \code{NULL} \\
\code{base\_date} & the start date as an integer vector \code{c(year, month, day)}, or
  \code{NULL} if the header gives none \\
\code{sampfrom} & the number of the first sample returned \\
\code{comments} & a character vector of the info strings of the header (without the
  \code{\#}) \\
\code{checksum\_ok} & \code{FALSE} if the samples do not match the checksums of the
  header; checksums are checked only when a read reaches the end of the record, so a
  shorter read gives \code{TRUE} \\
\end{xltabular}

\begin{xltabular}{\linewidth}{@{}lL@{}}
\caption{The elements of the list of class \code{"wfdbx\_header"} returned by
  \code{read\_header}}
\label{tab:r-header} \\
\toprule
\tabhead Element & \tabhead Contents \\
\midrule
\endfirsthead
\toprule
\tabhead Element & \tabhead Contents \\
\midrule
\endhead
\bottomrule
\endfoot
\code{name} & the record name, as given to \code{read\_header} \\
\code{fs} & the sampling frequency in Hz, or 0 if the header gives none \\
\code{counter\_freq} & the counter frequency, as \code{getcfreq()} returns it (the
  sampling frequency if the header gives none) \\
\code{base\_counter} & the base counter value, as \code{getbasecount()} returns it \\
\code{nsamples} & the number of samples per signal, or 0 if the header does not give it \\
\code{nsig} & the number of signals \\
\code{siginfo} & a list of \code{siginfo} lists, one per signal; \code{desc} has no white
  space at either end, and \code{adcres} is 0 if the header does not give it \\
\code{base\_time}, \code{base\_time\_ms}, \code{base\_date} & the start time and date, as
  in the result of \code{read\_record} \\
\code{comments} & a character vector of the info strings of the header (without the
  \code{\#}) \\
\code{metadata} & the metadata entries of the header, as \code{read\_metadata} returns
  them: a named character vector, in file order (empty for an EDF file and for the
  standard input) \\
\end{xltabular}

\begin{table}[htbp]
\centering
\caption{The columns of the data frame returned by \code{read\_annotations}, one row per
  annotation, in file order}
\label{tab:r-annotations}
\begin{tabularx}{\linewidth}{@{}lL@{}}
\toprule
\tabhead Column or attribute & \tabhead Contents \\
\midrule
\code{sample} & the annotation times, in samples \\
\code{symbol} & the mnemonics of the codes, such as \code{"N"}, \code{"V"} or \code{"+"}
  (as \code{annstr} gives them) \\
\code{code} & the annotation codes, such as \code{NORMAL} or \code{PVC} \\
\code{subtype} & the subtype of each annotation \\
\code{chan} & the channel number of each annotation \\
\code{num} & the annotator number of each annotation \\
\code{aux} & the auxiliary texts of the annotations, with \code{NA} where there is
  none \\
\code{attr(, "record")} & the record name, as given \\
\code{attr(, "annotator")} & the annotator name, such as \code{"atr"} \\
\code{attr(, "fs")} & the frequency in which \code{sample} is counted, normally the
  sampling frequency of the record; 0 if neither the record, nor the file, nor the
  argument \code{fs} of \code{read\_annotations} gives one \\
\bottomrule
\end{tabularx}
\end{table}

\begin{table}[htbp]
\centering
\caption{The list of class \code{"wfdbx\_check\_report"} returned by \code{check\_record},
  and its methods}
\label{tab:r-checkreport}
\begin{tabularx}{\linewidth}{@{}lL@{}}
\toprule
\tabhead Element or method & \tabhead Contents \\
\midrule
\code{record} & the record name, as given \\
\code{ok} & \code{TRUE} if no issue has the severity \code{"error"} \\
\code{issues} & a data frame with one row per problem, in the order of the checks, and
  the columns \code{severity} (\code{"error"}, \code{"warning"} or \code{"info"}),
  \code{code} (a stable identifier of the kind of problem, such as \code{"checksum"} or
  \code{"gain\_missing"}, \cref{tab:python-issue-codes}), \code{where} (\code{""} for the
  record, \code{"signal N"}, \code{"segment SEG"}, \code{"segment SEG, signal N"} or
  \code{"annotator ANN"}) and \code{message} \\
\code{format(report)} & the text form of the report, as a character vector: one line per
  issue, \code{"SEVERITY: WHERE: MESSAGE"} (\code{"SEVERITY: MESSAGE"} when \code{where}
  is empty), or the one line \code{no problems found} \\
\code{print(report)} & prints the lines of \code{format(report)}, and returns the report
  invisibly \\
\bottomrule
\end{tabularx}
\end{table}

\begin{xltabular}{\linewidth}{@{}lL@{}}
\caption{The elements of a \code{siginfo} list, the specification of a signal}
\label{tab:r-siginfo} \\
\toprule
\tabhead Element & \tabhead Contents \\
\midrule
\endfirsthead
\toprule
\tabhead Element & \tabhead Contents \\
\midrule
\endhead
\bottomrule
\endfoot
\code{fname} & the name of the signal file, as the header gives it \\
\code{desc} & the signal description, such as \code{"MLII"}, or \code{NULL} \\
\code{units} & the physical units, such as \code{"mV"}, or \code{NULL} if not specified \\
\code{gain} & the gain, in ADC units per physical unit (0: uncalibrated) \\
\code{initval} & the value of the first sample \\
\code{group} & the signal group number (the signals of one group share one file) \\
\code{fmt} & the storage format, such as \code{16}, \code{212} or \code{516} \\
\code{spf} & the number of samples per frame (more than 1 for the faster signals of a
  multi-frequency record) \\
\code{bsize} & the block size, in bytes (0 for ordinary files) \\
\code{adcres} & the ADC resolution, in bits (0: not specified) \\
\code{adczero} & the ADC output for an input of 0\,V \\
\code{baseline} & the ADC value that corresponds to 0 physical units \\
\code{nsamp} & the number of samples (0 if unknown) \\
\code{cksum} & the 16-bit checksum of the samples \\
\end{xltabular}

% ------------------------------------------------------------------------------------------
\subsection{Warnings and errors}
\label{sec:r-warnings}

\paragraph{Errors.} The high-level functions report errors as R errors (conditions of
class \code{"error"}, raised by \code{stop}), whose messages are given in the entries of
this section; \code{tryCatch} handles them. Each function checks its arguments before it
creates or changes a file, so an invalid argument leaves the files as they were. The
library may also print a message on the standard error when it cannot open a file (as in
\code{init: can't open header for record REC}); \code{wfdbquiet()} turns these messages
off (\cref{sec:r-library}).

The errors of the metadata functions, and those of the argument \code{metadata} of
\code{write\_record}, also have a class that tells their kind (\cref{tab:r-errors}), so
that a program can handle each kind on its own, with
\code{tryCatch(..., wfdbx\_file\_error = function(e) ...)}.

\begin{table}[htbp]
\centering
\caption{The classes of the errors of the metadata functions}
\label{tab:r-errors}
\begin{tabularx}{\linewidth}{@{}lL@{}}
\toprule
\tabhead Class & \tabhead Raised when \\
\midrule
\code{wfdbx\_file\_error} & a header cannot be found, read or written, or the directory
  given as \code{dir} does not exist (the errors that Python reports as \code{IOError}) \\
\code{wfdbx\_argument\_error} & an argument is invalid (a record name that is empty or
  \code{"-"}, a key or value that is not valid or not a string, changes or metadata that
  are not a named list or named character vector, two keys that differ only in letter
  case), or the header does not allow the change (more than one entry for the key, no
  record line, a header read from a URL without \code{dir}); Python's \code{ValueError}
  and \code{TypeError} \\
\bottomrule
\end{tabularx}
\end{table}
Both classes are also of class \code{"error"} (and \code{"condition"}), so
\code{tryCatch(..., error = )} catches them too.

\paragraph{The checksum warning.} When \code{read\_record} reads to the end of a record
whose samples do not match the checksums of its header, it returns all the samples and
issues an R warning: ``checksum error in record REC: the signal data do not match the
checksums in its header''. \code{checksum\_warning = FALSE} turns the warning off;
\code{checksum\_ok} of the result is \code{FALSE} either way.

\paragraph{The duplicate-key warning.} When a header holds two metadata entries whose keys
differ only in letter case (written by other software), \code{read\_metadata},
\code{get\_metadata} and \code{read\_header} use the first and issue an R warning:
``header of record REC has more than one entry for key KEY (lines L1, L2); the first is
used''. \code{set\_metadata} and \code{update\_metadata} then refuse to change that key
(an error of class \code{wfdbx\_argument\_error}), and \code{remove\_metadata} removes all
of its entries.

Both warnings are issued with \code{warning()}, without the call. As any other R warning,
they can be handled (\code{withCallingHandlers}), silenced (\code{suppressWarnings}) or
turned into errors (\code{options(warn = 2)}):

\begin{lstlisting}[language=R,style=example]
library(wfdbx)

# A record whose last sample is changed after it was written: the checksum fails.
write_record("ck", 0:99, fs = 100, gains = 100, adcres = 12, units = "mV")
con <- file("ck.dat", "r+b")
invisible(seek(con, 198, rw = "write"))
writeBin(as.raw(c(0, 0)), con)
close(con)
w <- NULL
rec <- withCallingHandlers(read_record("ck"), warning = function(c) {
  w <<- conditionMessage(c)
  invokeRestart("muffleWarning")
})
cat(dim(rec$signals), rec$checksum_ok, "\n")    # 100 1 FALSE: all the samples
print(w)
# "checksum error in record ck: the signal data do not match the checksums in its header"
rec <- read_record("ck", checksum_warning = FALSE)      # no warning

# Two entries for one key: the first is used, with a warning.
write_record("dup", rep(0, 10), fs = 100, gains = 0, adcres = 12,
             comments = c(" Age: 63", " AGE: 64"))
print(suppressWarnings(get_metadata("dup", "age")))     # "63"
tryCatch(set_metadata("dup", "Age", "65"),
         wfdbx_argument_error = function(e) print(conditionMessage(e)))
# "header of record dup has more than one entry for key Age (lines 3, 4); remove them first"
print(remove_metadata("dup", "age"))                    # 2
\end{lstlisting}

\begin{lstlisting}[language=R,style=example]
library(wfdbx)

tryCatch(read_record("nosuchrecord"), error = function(e) print(conditionMessage(e)))
# "can't read record nosuchrecord" (after the library's message on the standard error)

write_record("raw", 0:9, fs = 100, gains = 0, adcres = 12)      # uncalibrated
tryCatch(read_record("raw", physical = TRUE), error = function(e) print(conditionMessage(e)))
# "signal 0 of record raw is uncalibrated"

# The metadata functions: the class tells the kind of error.
e <- tryCatch(read_metadata("nosuchrecord"), error = function(e) e)
print(class(e))           # "wfdbx_file_error" "error" "condition"
e <- tryCatch(set_metadata("raw", "a:b", "1"), error = function(e) e)
print(class(e))           # "wfdbx_argument_error" "error" "condition"
print(conditionMessage(e))   # "metadata key 'a:b' must not contain ':'"
\end{lstlisting}

% ==========================================================================================
\section{The WFDB library functions in R}
\label{sec:r-library}

Every function of the WFDB library is available under its C name, with the same
arguments, in the same order, and the same meaning, so that programs and documentation
written for the C library carry over directly. The functions, their arguments and their
results are documented once, for all the languages, in \cref{chap:library}; this section
gives the rules by which the C declarations become R functions. The help pages
\code{?isigopen}, \code{?timstr}, \code{?annopen}, \code{?annstr}, \code{?calopen},
\code{?getwfdb}, \code{?ecgcodes}, \code{?isann} and \code{?siginfo} list the functions of
each part of the library, with the result of each.

\subsection{Conventions}

\paragraph{Names.} The functions, the constants (\code{WFDB\_READ},
\code{WFDB\_INVALID\_SAMPLE}, \code{WFDB\_HIGHRES}, \ldots) and the annotation codes
(\code{NORMAL}, \code{PVC}, \ldots, \cref{chap:formats}; \code{ecgcodes()} returns them as
a list) keep their C names. The macros of \file{<wfdb/ecgmap.h>} are functions, under both
their macro names and the names of the functions behind them: \code{isqrs} and
\code{wfdb\_isqrs}, \code{map1} and \code{wfdb\_map1}, and so on; they take vectors. The C
function \code{sample} is \code{wfdb\_sample}, so that it does not mask R's
\code{sample}.

\paragraph{Numbers} are doubles. Signal, group and annotator numbers start at 0, as in C
(signal \code{s} is element \code{s + 1} of an R vector). C's integer arithmetic (wrap
around, casts to \code{short} or \code{char}) is applied where C applies it.

\paragraph{NULL is \code{NULL}.} A \code{NULL} pointer argument or result is
\code{NULL}: \code{sampfreq(NULL)} returns the frequency of the open record,
\code{getinfo(NULL)} continues with the info strings of the last record, and
\code{getinfo} returns \code{NULL} after the last one. A list element that is
\code{NULL} is set with \code{x["f"] <- list(NULL)} (\code{x\$f <- NULL} removes it).

\paragraph{Structures} are lists with the C field names, made by the constructors
\code{siginfo()}, \code{anninfo()}, \code{annotation()}, \code{calinfo()} and
\code{seginfo()} (\cref{sec:r-structures}). An array of structures is a list of them:
\code{annopen(record, list(anninfo("atr", WFDB\_READ)))}, \code{osigfopen(list(siginfo(...),
...))}.

\paragraph{Arrays and output arguments.} R functions do not change their arguments: where
a C function fills an array or a structure that the caller passes, the R function returns
it instead, in a list with the C return value:
\begin{lstlisting}[language=R]
r <- isigopen(record, nsig)      # r$nsig: the C result; r$si: the siginfo lists
r <- osigopen(record, nsig)      # the same
g <- getvec()                    # g$stat: the C result; g$v: one sample per signal
g <- getframe()                  # g$stat; g$v: one sample per sample of a frame
g <- getann(an)                  # g$stat; g$annot: the annotation (NULL on errors)
g <- getcal(desc, units)         # g$stat; g$cal: the calinfo list (NULL if none)
r <- wfdbinit(record, aiarray, nsig)   # r$stat; r$si
s <- getseginfo()                # a list of seginfo lists (empty: not multi-segment)
\end{lstlisting}
With \code{nsig} $\le 0$, \code{isigopen} reads the header only: \code{r\$nsig} is then
the number of signals of the record, and \code{r\$si} the specifications of up to
\code{-nsig} of them. The arguments that C takes only to receive results are left out
(\code{getvec()} has none), and the lengths of arrays have defaults: \code{osigfopen(siarray,
nsig = length(siarray))}, \code{setheader(record, siarray, nsig = length(siarray))},
\code{setmsheader(record, segment\_name, nsegments = length(segment\_name))},
\code{annopen(record, aiarray, nann = NULL)}. \code{putvec(vector)}, \code{putann(an,
annot)}, \code{ungetann(an, annot)} and \code{putcal(cal)} take the values to write.

\paragraph{Strings} are character strings, handled as bytes (ISO-8859-1, never
re-encoded), so every byte sequence round-trips unchanged; functions that take a string
also take a raw vector. The \code{aux} element of an annotation is \code{NULL}, or a raw
vector whose first byte is the length of the text that follows, as in C;
\code{aux\_string(aux)} returns the text (up to its first NUL byte), or \code{NULL}.

\paragraph{Return values and errors.} The functions return what the C functions return:
a status (negative on errors), a number of signals, a frequency, a string or
\code{NULL}. Errors are reported by these values and by a message on the standard error,
as in C, never by R errors. \code{wfdberror()} returns the last message;
\code{wfdbquiet()} turns the messages off and \code{wfdbverbose()} on again;
\code{wfdb\_error(fmt, ...)} reports a message of a program's own in the same way.

\paragraph{State.} As in C, the library keeps its state in the R session: the open record,
its input signals, the output signals and the annotators remain open between calls,
until \code{wfdbquit()} closes them. \code{wfdbquit()} leaves the input and output
directories and the network setting as they are. The high-level functions of
\cref{sec:r-api} call \code{wfdbquit()}, which closes anything opened with the library
functions.

\paragraph{Additions.} \code{getframes(n, minstat = 0)} reads up to \code{n} sample vectors
at once (\code{n} may be \code{Inf}), as \code{n} calls of \code{getvec} would, but
decoding whole blocks of the signal files. It returns \code{list(v, stat)}: a matrix with
one row per vector read and one column per signal, and the \code{getvec} status of the
first vector that could not be read, or the number of signals if all \code{n} were read;
with \code{minstat}, it stops at the first vector for which \code{getvec} returns less
than \code{minstat}. \code{putframes(m)} writes the rows of a matrix, as \code{putvec} on
each row. \code{sample\_valid()} returns the validity flag of the last
\code{wfdb\_sample}.

\paragraph{Speed.} Reading sample by sample with \code{getvec} in an R loop is much slower
than reading in blocks: \code{read\_record} and \code{getframes} decode whole blocks with
vectorized code, and are the way to read large records (in chunks, with \code{sampfrom}
and \code{sampto}, for a record too large for memory).

\subsection{The structures}
\label{sec:r-structures}

The structures of the C library are lists whose elements have the names of the fields of
the C structures (\cref{tab:r-structures}). Their constructors take the elements as
arguments, in this order, all with defaults:
\begin{lstlisting}[language=R]
siginfo(fname = NULL, desc = NULL, units = NULL, gain = 0,
        initval = 0, group = 0, fmt = 0, spf = 1, bsize = 0,
        adcres = 0, adczero = 0, baseline = 0, nsamp = 0,
        cksum = 0)
anninfo(name = NULL, stat = WFDB_READ)
annotation(time = 0, anntyp = 0, subtyp = 0, chan = 0, num = 0,
           aux = NULL)
calinfo(low = 0, high = 0, scale = 0, sigtype = NULL,
        units = NULL, caltype = 0)
seginfo(recname = "", nsamp = 0, samp0 = 0)
\end{lstlisting}
For a list of \code{n} structures, \code{replicate(n, siginfo(), simplify = FALSE)}.

\begin{xltabular}{\linewidth}{@{}lL@{}}
\caption{The lists of the C structures}
\label{tab:r-structures} \\
\toprule
\tabhead Constructor (C structure) & \tabhead Elements \\
\midrule
\endfirsthead
\toprule
\tabhead Constructor (C structure) & \tabhead Elements \\
\midrule
\endhead
\bottomrule
\endfoot
\code{siginfo} (\code{WFDB\_Siginfo}) & \code{fname}, \code{desc}, \code{units},
  \code{gain}, \code{initval}, \code{group}, \code{fmt}, \code{spf}, \code{bsize},
  \code{adcres}, \code{adczero}, \code{baseline}, \code{nsamp}, \code{cksum}
  (\cref{tab:r-siginfo}) \\
\code{seginfo} (\code{WFDB\_Seginfo}) & \code{recname} (the name of the segment),
  \code{nsamp} (its number of samples), \code{samp0} (the number of its first sample in
  the record) \\
\code{anninfo} (\code{WFDB\_Anninfo}) & \code{name} (the annotator name, the suffix of the
  file), \code{stat} (\code{WFDB\_READ} or \code{WFDB\_WRITE} for the MIT format,
  \code{WFDB\_AHA\_READ} or \code{WFDB\_AHA\_WRITE} for the AHA format) \\
\code{annotation} (\code{WFDB\_Annotation}) & \code{time} (in sample intervals),
  \code{anntyp} (the annotation code), \code{subtyp}, \code{chan}, \code{num},
  \code{aux} (\code{NULL}, or a raw vector whose first byte is the length of the text
  that follows) \\
\code{calinfo} (\code{WFDB\_Calinfo}) & \code{low}, \code{high} (the levels of the
  calibration pulse, in physical units), \code{scale} (physical units per cm),
  \code{sigtype} (matched against the beginning of a signal description), \code{units},
  \code{caltype} (\code{WFDB\_AC\_COUPLED} or \code{WFDB\_DC\_COUPLED}, or-ed with
  \code{WFDB\_CAL\_SQUARE}, \code{WFDB\_CAL\_SINE} or \code{WFDB\_CAL\_SAWTOOTH}) \\
\end{xltabular}

\subsection{The settings functions}
\label{sec:r-settings}

The functions of \cref{tab:r-settings} are additions of WFDB-x to the library: they set
where files are read and written, network access and legacy mode, for the high-level
functions, the library functions and the tools alike (\cref{chap:files,chap:behavior}).
Each \code{set} function returns the previous setting, invisibly, so that it can be
restored. The C function \code{setwfdb(p)} sets the input directory too, and
\code{getwfdb()} returns it; in legacy mode they set and return the WFDB path, as in C.

\begin{xltabular}{\linewidth}{@{}lL@{}}
\caption{The settings functions}
\label{tab:r-settings} \\
\toprule
\tabhead Function & \tabhead Effect and result \\
\midrule
\endfirsthead
\toprule
\tabhead Function & \tabhead Effect and result \\
\midrule
\endhead
\bottomrule
\endfoot
\code{wfdbx\_set\_input\_dir(dir)} & sets the input directory: a local directory, a
  \code{\textasciitilde} path or a URL; \code{NULL} or \code{""} restores the default
  (\code{WFDBX\_INPUT\_DIR}, or else the current directory). Returns the previous input
  directory, or \code{NULL} if none was set \\
\code{wfdbx\_get\_input\_dir()} & returns the input directory, or \code{NULL} if none is
  set \\
\code{wfdbx\_set\_output\_dir(dir)} & sets the output directory: a local directory or a
  \code{\textasciitilde} path (a URL is refused, with a message, and the setting is
  unchanged); \code{NULL} or \code{""} restores the default (\code{WFDBX\_OUTPUT\_DIR}, or
  else the current directory). Returns the previous output directory, or \code{NULL} \\
\code{wfdbx\_get\_output\_dir()} & returns the output directory, or \code{NULL} if none
  is set \\
\code{wfdbx\_set\_network(on)} & turns network access on (\code{TRUE} or 1) or off
  (\code{FALSE} or 0); \code{NULL} or a negative number restores the default (on if
  \code{WFDBX\_NETWORK} is a nonzero number). Returns the previous setting, 1 or 0 \\
\code{wfdbx\_get\_network()} & returns 1 if network access is on, 0 otherwise \\
\code{wfdbx\_set\_legacy(on)} & turns legacy mode on (\code{TRUE}) or off, for the whole
  library. Returns the previous setting, \code{TRUE} or \code{FALSE} \\
\code{wfdbx\_get\_legacy()} & returns \code{TRUE} if legacy mode is on, \code{FALSE}
  otherwise \\
\code{wfdbx\_legacy\_args(argv)} & if \code{argv[2]} (\code{argv}: the argument vector,
  program name first) is \code{"--legacy"}, turns legacy mode on; returns
  \code{list(argv, legacy)}: the arguments without \code{"--legacy"}, and whether it was
  found \\
\code{wfdbx\_time\_needs\_freq(string)} & returns 1 if \code{strtim(string)} needs a
  sampling frequency, 0 otherwise \\
\code{wfdbx\_write\_errors()} & returns the number of write errors since the R session
  started (output files that could not be written, flushed or closed) \\
\code{wfdbx\_note\_write\_error()} & counts a write error, for a program that writes a
  file itself \\
\code{wfdbx\_input\_path(name)} & returns the path where a program finds a file that it
  reads itself, as \code{wfdb\_open} finds a file of that name: in normal mode, a
  relative name under the input directory, if one is set, a \code{\textasciitilde} name
  expanded, an absolute name as given; in legacy mode, the name as given. \code{""} and
  \code{"-"} are returned as given; \code{NULL} for \code{NULL} or a URL, or when the
  input directory makes the name a URL \\
\code{wfdbx\_output\_path(name)} & returns the path where a program creates a file that
  it writes itself, as \code{wfdb\_open} creates a file of that name: in normal mode, a
  relative name under the output directory, if one is set, a \code{\textasciitilde} name
  expanded, an absolute name as given; in legacy mode, any name under the output
  directory, if one is set, as a plain prefix. When the output directory is prefixed,
  the missing directories of the path are created. \code{""} and \code{"-"} are returned
  as given; \code{NULL} for \code{NULL} or a URL \\
\end{xltabular}

\subsection{Examples}

Reading the samples and the annotations of a record, as a C program does (the record
\code{100s} is read from the current directory, a copy of \file{tests/data}):

\begin{lstlisting}[language=R,style=example]
library(wfdbx)

r <- isigopen("100s", 2)                 # open 2 signals
nsig <- r$nsig
cat(nsig, sapply(r$si, function(s) paste(s$desc, s$fmt, s$gain)), "\n")
# 2 MLII 212 200 V5 212 200
print(sampfreq(NULL))                    # 360: the frequency of the open record

stat <- isigsettime(strtim("0:30"))      # skip to 30 seconds
g <- getvec()                            # list(stat, v)
cat(g$stat, g$v, sapply(0:(nsig - 1), function(k) aduphys(k, g$v[k + 1])), "\n")
# 2 947 968 -0.385 -0.28

f <- getframes(360)                      # the next second, as a matrix
cat(dim(f$v), f$stat, "\n")              # 360 2 2

stat <- annopen("100s", list(anninfo("atr", WFDB_READ)))
n <- 0
while ((g <- getann(0))$stat == 0) {     # 0: an annotation was read; -1: the end
  if (n < 2) cat(mstimstr(g$annot$time), annstr(g$annot$anntyp),
                 aux_string(g$annot$aux), "\n")
  n <- n + 1
}
cat(n, "annotations\n")                  # 75 annotations

wfdbquiet()                              # no messages on the standard error
cat(isigopen("nosuchrecord", 0)$nsig, wfdberror())
# -1 init: can't open header for record nosuchrecord
wfdbverbose()
stat <- wfdbquit()                       # close everything
\end{lstlisting}

Writing a record and an annotation file, sample by sample and annotation by annotation:

\begin{lstlisting}[language=R,style=example]
library(wfdbx)

out <- list(siginfo(fname = "ramp.dat", desc = "ramp", units = "mV", gain = 200,
                    fmt = 16, adcres = 12))
print(osigfopen(out))                    # 1: one output signal
stat <- setsampfreq(250)
for (k in 0:499) putvec(k)
print(newheader("ramp"))                 # 0: ramp.hea is written
stat <- annopen("ramp", list(anninfo("qrs", WFDB_WRITE)))
for (t in c(100, 350)) putann(0, annotation(time = t, anntyp = NORMAL))
stat <- wfdbquit()                       # completes ramp.dat and ramp.qrs

print(readLines("ramp.hea", n = 1))              # "ramp 1 250 500"
print(read_annotations("ramp", "qrs")$sample)    # 100 350
\end{lstlisting}

% ==========================================================================================
\section{The command-line tools in R}
\label{sec:r-tools}

The 33 command-line tools (\code{rdsamp}, \code{rdann}, \code{wrsamp}, \code{xform},
\ldots) are run by the driver script \file{wfdbx.R} with \cmd{Rscript}, followed by the
name of the tool and its arguments. The tools, their options, their output and their exit
status are those of the C tools, and are documented in \cref{chap:tools}. In a copy of
the repository the script is \file{r/inst/bin/wfdbx.R}, which runs without the package
being installed; in an installed package its path is \code{system.file("bin", "wfdbx.R",
package = "wfdbx")}. A shell alias makes it convenient:

\begin{lstlisting}[style=shell]
alias wfdbx='Rscript /path/to/wfdb-x/r/inst/bin/wfdbx.R'
export WFDBX_INPUT_DIR=/path/to/wfdb-x/tests/data   # the directory of 100s
wfdbx rdsamp -r 100s -t 10 -p     # the first 10 seconds of 100s, in physical units
wfdbx rdann -r 100s -a atr        # its annotations
wfdbx rdsamp --legacy -r 100s     # legacy mode: --legacy as the first argument
\end{lstlisting}

The driver script also takes these options, instead of a tool:

\begin{table}[htbp]
\centering
\caption{The options of the driver script \file{wfdbx.R}}
\label{tab:r-wfdbx-options}
\begin{tabularx}{\linewidth}{@{}lL@{}}
\toprule
\tabhead Command & \tabhead Effect \\
\midrule
\cmd{Rscript wfdbx.R TOOL [ARGUMENTS ...]} & runs the tool \code{TOOL}, and exits with
  its status \\
\cmd{Rscript wfdbx.R --list} & lists the names of the tools, one per line \\
\cmd{Rscript wfdbx.R --version} & prints the version, \code{WFDB-x 1.2.1} \\
\cmd{Rscript wfdbx.R --help}, \cmd{Rscript wfdbx.R -h} & prints the usage \\
\bottomrule
\end{tabularx}
\end{table}

Without arguments, the script prints the usage and exits with status 1; an unknown tool
is an error (``wfdbx: unknown tool 'TOOL' (try 'wfdbx -{}-list')'', status 1), and an R
error in a tool is reported as \code{wfdbx: error: MESSAGE} (status 1; with the
environment variable \code{WFDBX\_DEBUG} set, the call stack follows). When an output file
or the standard output cannot be written (a full disk, for example), a tool reports it and
exits with status 1 (legacy mode keeps the original status, 0). The files that a tool
reads or creates itself follow the input and output directories, as the records do: the
\opt{-i} file of \code{wrsamp}, \code{edf2mit}, \code{wav2mit}, \code{a2m} and
\code{ad2m}, the input files of \code{ahaecg2mit} and the \opt{-S} script of
\code{xform} are read from the input directory, and the WAV file of \code{mit2wav}, the
EDF file of \code{mit2edf}, the master header of \code{snip -m} and the \file{.des},
\file{.ecg}, \file{.key} and \file{.txt} files of \code{parsescp} are created in the
output directory (\code{wfdbx\_input\_path} and \code{wfdbx\_output\_path},
\cref{tab:r-settings}); such a file cannot be read from a URL. The script runs the tools
in the C locale, with R's warnings turned off, so that their output is that of the C
tools.

\paragraph{The start-up cache.} When run from a copy of the repository, the script keeps a
compiled copy of the sources in \file{\textasciitilde/.cache/R/wfdbx/} (or under
\code{\$XDG\_CACHE\_HOME}), so that a tool starts quickly; \code{WFDBX\_NO\_CACHE=1} turns
this off.

\paragraph{Windows.} The tools write their output in binary through \file{/dev/stdout} and
\file{/dev/stderr}. On Windows these do not exist, so the output is written as text and
NUL bytes are lost.

\paragraph{The tools from R.} \code{wfdbx\_tool} runs a tool in the R session, and
\code{wfdbx\_tools} lists them.

\apientry{wfdbx\_tool}{Run a command-line tool}
\begin{lstlisting}[language=R,style=signature]
wfdbx_tool(tool, args = character(0))
\end{lstlisting}
Runs one of the command-line tools with the given arguments, as the command line would,
in the R session. The output of the tool goes directly to the standard output and error
of the R process: \code{sink} and \code{capture.output} do not capture it, and a graphical
console may not show it. To capture the output, the tool is run in a separate process
with the driver script, as in \code{system2("Rscript", c(driver, "rdsamp", ...), stdout =
TRUE)}.
\begin{params}
\item[tool] The name of the tool, a character string (one of \code{wfdbx\_tools()}).
\item[args] The command-line arguments of the tool, as a character vector (one element
  per argument, without the program name). Default \code{character(0)}: none.
\end{params}
\returns{The exit status of the tool (0 on success), a number.}
\raises{Stops with an error if \code{tool} is not the name of a tool (``unknown tool
'TOOL'\,'').}
\begin{lstlisting}[language=R,style=example]
library(wfdbx)

status <- wfdbx_tool("rdsamp", c("-r", "100s", "-t", "s2"))     # the first 2 samples
# rdsamp prints two lines: the sample numbers 0 and 1, and the samples 995 1011
print(status)                            # 0
\end{lstlisting}

\apientry{wfdbx\_tools}{List the command-line tools}
\begin{lstlisting}[language=R,style=signature]
wfdbx_tools()
\end{lstlisting}
Returns the names of the command-line tools.
The function has no parameters.
\returns{A character vector of the names of the tools, in alphabetical order (C
collation).}
\raises{Nothing.}
\begin{lstlisting}[language=R,style=example]
library(wfdbx)

print(length(wfdbx_tools()))             # 33
print(head(wfdbx_tools(), 4))            # "a2m" "ad2m" "ahaecg2mit" "ann2rr"
\end{lstlisting}

\paragraph{Environment variables.} The tools, like the rest of the package, use the
environment variables of \cref{tab:r-environment}. They are read when they are first
needed in an R session.

\begin{table}[htbp]
\centering
\caption{The environment variables used by the R package}
\label{tab:r-environment}
\begin{tabularx}{\linewidth}{@{}lL@{}}
\toprule
\tabhead Variable & \tabhead Meaning \\
\midrule
\code{WFDBX\_INPUT\_DIR} & the input directory (a directory, a \code{\textasciitilde}
  path or a URL), when \code{wfdbx\_set\_input\_dir} does not set one \\
\code{WFDBX\_OUTPUT\_DIR} & the output directory, when \code{wfdbx\_set\_output\_dir} does
  not set one \\
\code{WFDBX\_NETWORK} & network access is on if it is a nonzero number, when
  \code{wfdbx\_set\_network} does not set it \\
\code{WFDBANNSORT} & \code{0} to keep annotation files in the order written (no sorting
  when they are closed) \\
\code{WFDBGVMODE} & the default getvec mode (\code{setgvmode} with a negative mode) \\
\code{WFDBCAL} & the calibration file read by \code{calopen(NULL)} (default
  \file{wfdbcal}) \\
\code{WFDB} & the WFDB path, used in legacy mode only (outside legacy mode, it is
  ignored, with a message) \\
\code{WFDB\_FLAC\_CHECK\_MD5} & if set, the MD5 signature of FLAC signal files is
  verified when they are read \\
\code{WFDBX\_DEBUG} & if set, the driver script prints the call stack of an R error in a
  tool \\
\code{WFDBX\_NO\_CACHE} & \code{1} to turn the start-up cache of the driver script off \\
\code{XDG\_CACHE\_HOME} & the directory of the start-up cache (default
  \file{\textasciitilde/.cache}) \\
\bottomrule
\end{tabularx}
\end{table}

Remote files are downloaded with the program \cmd{curl} if it is installed, and with R's
\code{url()} otherwise; each is read once and kept for the rest of the R session, and a
download gives up when it cannot connect, or receives nothing, for 30 seconds.

% ==========================================================================================
\section{Testing}
\label{sec:r-testing}

The unit tests of the R package are in \file{r/tests}: the runner \file{run\_tests.R} and
the files \file{unit/test\_*.R}, written with base R only. From the root of the
repository:

\begin{lstlisting}[style=shell]
Rscript r/tests/run_tests.R              # all the unit tests
Rscript r/tests/run_tests.R metadata     # the files whose name matches "metadata"
\end{lstlisting}

The runner loads the sources of \file{r/R} (or the installed package, under \cmd{R CMD
check}), runs each test file, and prints one line per file and one per failed check; it
exits with status 1 if a check fails. The tests check the library functions with expected
values of the C library, the high-level functions, the bulk paths against the
sample-by-sample paths, \code{printf} and \code{strtod} against the GNU C library, the
FLAC codec, \code{check\_record}, and the header metadata functions against the cases
shared by the seven languages (in \file{tests/metadata}: results, warnings and the
headers, byte for byte). Network access is tested against a local HTTP server
(127.0.0.1, run with \cmd{python3}; skipped without it). The tests work in temporary
directories, with copies of the test records of \file{tests/data}; the tests that need
\file{tests/data} are skipped when it is not available (in the installed package).

The package check builds the package and runs its examples and tests:

\begin{lstlisting}[style=shell]
R CMD build r && R CMD check --no-manual wfdbx_*.tar.gz
\end{lstlisting}

The cross-language conformance suite compares the output of the tools with that of the
original C programs, case by case: 1135 cases, and 883 in legacy mode, which runs each
tool with \opt{--legacy}. It runs the driver script \file{r/inst/bin/wfdbx.R}. From the
root of the repository:

\begin{lstlisting}[style=shell]
PYTHONPATH=$PWD/python/src python3 tests/conformance/harness.py check --impl r -j8
PYTHONPATH=$PWD/python/src python3 tests/conformance/harness.py check --impl r -j8 --legacy
\end{lstlisting}
